\documentclass[12pt]{article}
\usepackage{amsmath,amsfonts,fullpage,amsthm,graphicx}

\newtheorem*{theorem}{Theorem}
\newtheorem*{fact}{Fact}
\newtheorem*{goal}{Goal}
\newtheorem*{idea}{Idea}
\newtheorem*{defn}{Definition}
\newtheorem*{ex}{Example}

%Style section
%\setlength{\textheight}{10in}
%\setlength{\textwidth}{7in}
%\hoffset=-0.75in
 \pagestyle{plain}
% \setlength{\topmargin}{0in}
% \setlength{\headsep}{0in}
% \setlength{\oddsidemargin}{0in}
% \setlength{\evensidemargin}{0in}

\begin{document}


\pagestyle{empty} %use this to get rid of page numbers

\noindent
{Math 166 
%- Dr.\ Jessica Striker
 \hskip 1.5 truein  Names:\ \hrulefill}

\noindent
%\vskip 0.3truein
%\bigskip
%\noindent
{Series convergence}

\noindent  \hskip 3.5 truein  Section Number:\ \hrulefill
\medskip

\begin{goal} \rm
Explore more methods of determining whether series converge or diverge. 
\end{goal}
%\smallskip

\begin{enumerate}
\item State the integral test in your own words. What three conditions do you need to check in order to use the integral test?

%\vspace{.5in}
%\item In your team, each person should attempt to use \textbf{one} of the following tests to determine whether the series  converges or diverges: $n$th term divergence test, Comparison test, Limit comparison test, Integral test. Which test(s) worked well for this series? 

\vspace{.7in}
\item Use either the $n$th term divergence test or the integral test to determine whether the following series converge or diverge. Be sure to verify the hypotheses, state your conclusion, and say which test you used. 
%You should approach each problem together as a team, but you may each try different tests to more quickly determine what test will work in each problem.
\begin{enumerate}
\item $\displaystyle\sum_{n=1}^{\infty} (1.01)^n$

%\vspace{3.5in}
%\pagebreak
\vspace{2.5in}

\item $\displaystyle\sum_{n=1}^{\infty} e^{-3n}$

%\vspace{1.5in}
%\item $\displaystyle\sum_{n=2}^{\infty} \frac{1}{\sqrt{n^4+n^3}}$

%\vspace{1.5in}
%\item $\displaystyle\sum_{n=1}^{\infty} \frac{n^3-5n+1}{7n^5+2}$
\end{enumerate} 
\newpage
\item Use the integral test to determine for which values of $p$ the following series converges.
\[
\sum_{n = 2}^\infty \frac{1}{n (\ln n)^p}
\]
\vfill
%\item Evaluate the following series:
%\[
%\displaystyle \sum_{n = 4}^\infty \frac{2^n}{4^n} + \frac{1}{(n + 1)(n + 3)}.
%\]
%\vfill

%\newpage

\item For which integer values of $a > 0$ does the following series converge?  Explain.
\[
\sum_{n = a}^\infty \frac{1}{n}
\]
\vfill
\item For which values of $p$ does the following series converge?  Explain.
\[
\sum_{n = 5}^\infty \frac{5}{n^{2p + 9}}
\]
\vfill
\end{enumerate}
%\newpage
%
%\item State the {\bf integral test} in your own words. What three conditions do you need to check in order to use the integral test?
%
%%\vspace{.5in}
%%\item In your team, each person should attempt to use \textbf{one} of the following tests to determine whether the series  converges or diverges: $n$th term divergence test, Comparison test, Limit comparison test, Integral test. Which test(s) worked well for this series? 
%
%\vspace{.7in}
%\item Use either the $n$th term divergence test or the integral test to determine whether the following series converge or diverge. Be sure to verify the hypotheses, state your conclusion, and say which test you used. 
%%You should approach each problem together as a team, but you may each try different tests to more quickly determine what test will work in each problem.
%\begin{enumerate}
%\item $\displaystyle\sum_{n=1}^{\infty} (1.01)^n$
%
%%\vspace{3.5in}
%%\pagebreak
%\vspace{2.5in}
%
%%\item $\displaystyle\sum_{n=1}^{\infty} e^{-3n}$
%%
%%%\vspace{1.5in}
%%%\item $\displaystyle\sum_{n=2}^{\infty} \frac{1}{\sqrt{n^4+n^3}}$
%%
%%\vspace{2.5in}
%%\item $\displaystyle\sum_{n=1}^{\infty} \frac{7^n+1}{2^n}$
%%
%%\vspace{2.5in}
%\item $\displaystyle\sum_{n=2}^{\infty} \frac{n}{n^2+1}$
%
%\vspace{2.5in}
%\item $\displaystyle\sum_{n=1}^{\infty} \frac{\cos(n)}{n^2}$ (Hint: Can you use the alternating series test? Why not? Instead, check for absolute convergence.)
%\item $\displaystyle\sum_{n=1}^{\infty} \frac{1}{4+5^n}$
%\item $\displaystyle\sum_{n=3}^{\infty} \ln(n)$


%\end{enumerate}

%\end{enumerate}

\end{document}
